\chapter{Contoh Tugas, Portofolio, dan Rubrik Refleksi}
\label{app:contoh-tugas}

Lampiran ini menyediakan contoh instrumen yang dipakai sepanjang
semester: ringkasan tugas, kerangka portofolio, kerangka jurnal, dan
rubrik penilaian refleksi.

\section{Ringkasan Tugas Individu}

\begin{table}[H]
    \centering
    \caption{Ringkasan tujuh tugas individu dan Sub-CPMK yang
    diukurnya.}
    \label{tab:app-daftar-tugas}
    \footnotesize
    \begin{tabular}{@{}cp{3.0cm}p{5.4cm}p{2.2cm}@{}}
        \toprule
        \textbf{Tugas} & \textbf{Bab} & \textbf{Bentuk pekerjaan} &
        \textbf{Sub-CPMK} \\
        \midrule
        Tugas 1 & Bab~\ref{chap:register} &
        Laporan observasi register dan \eng{flag} dengan tabel
        pengamatan & 1.2 \\
        \addlinespace
        Tugas 2 & Bab~\ref{chap:format-instruksi} &
        Lembar kerja \eng{encoding} dan \eng{decoding} instruksi &
        2.1 \\
        \addlinespace
        Tugas 3 & Bab~\ref{chap:translasi} &
        Portofolio perakitan \file{.asm} $\rightarrow$
        \file{.obj} $\rightarrow$ \file{.exe} &
        2.3 \\
        \addlinespace
        Tugas 4 & Bab~\ref{chap:aritmatika} &
        Program kalkulator sederhana dengan penanganan
        \eng{overflow} & 3.1 \\
        \addlinespace
        Tugas 5 & Bab~\ref{chap:kontrol} &
        Program pencarian nilai maksimum dan pengurutan &
        3.3 \\
        \addlinespace
        Tugas 6 & Bab~\ref{chap:prosedur} &
        Prosedur berparameter melalui stack beserta gambar
        \eng{stack frame} & 3.4 \\
        \addlinespace
        Tugas 7 & Bab~\ref{chap:interupsi} &
        Program masukan/keluaran dengan \code{INT 21h} &
        2.4 \\
        \bottomrule
    \end{tabular}
\end{table}

\section{Contoh Lembar Tugas}

Format lembar tugas yang dikumpulkan mahasiswa adalah sebagai
berikut.

\begin{infobox}
\textbf{Lembar Tugas \dots}\\[0.3em]
\textbf{Nama / NIM:} \dots\\
\textbf{Bab:} \dots \quad \textbf{Sub-CPMK:} \dots\\
\textbf{Tujuan tugas:} \dots\\
\textbf{Langkah kerja:} (daftar bernomor)\\
\textbf{Kendala yang dihadapi:} \dots\\
\textbf{Cara mengatasi:} \dots\\
\textbf{Kesimpulan:} apa yang dipahami setelah mengerjakan tugas
ini\\
\textbf{Lampiran:} berkas \file{.asm}, tangkapan layar keluaran,
tabel pengamatan
\end{infobox}

\section{Kerangka Portofolio Program}

Portofolio disusun menurut bab dan memuat berkas berikut untuk
setiap bab.

\begin{table}[H]
    \centering
    \caption{Kerangka isi portofolio program.}
    \label{tab:app-portofolio}
    \small
    \begin{tabular}{@{}p{3.4cm}p{10.0cm}@{}}
        \toprule
        \textbf{Berkas} & \textbf{Isi} \\
        \midrule
        \file{babNN/} & Direktori per bab \\
        \addlinespace
        \file{*.asm} & Seluruh program yang ditulis, termasuk versi
        yang belum berhasil (diberi akhiran \file{\_gagal} pada
        komentar, bukan pada nama berkas) \\
        \addlinespace
        \file{catatan-babNN.md} & Catatan singkat: tujuan setiap
        program, kesalahan yang ditemukan, dan cara memperbaikinya \\
        \addlinespace
        \file{bukti/} & Tangkapan layar keluaran program dan tabel
        pengamatan register \\
        \bottomrule
    \end{tabular}
\end{table}

\begin{infobox}
Versi program yang gagal tetap disimpan. Justru versi gagal inilah
yang menunjukkan proses belajar Anda, dan menjadi bahan penilaian
\CPMK{04}. Menghapus versi gagal membuat riwayat perbaikan Anda
hilang.
\end{infobox}

\section{Kerangka Jurnal Refleksi Mingguan}

Setiap sesi praktikum ditutup dengan satu entri jurnal yang memuat
empat kalimat wajib berikut.

\begin{enumerate}
    \item \textbf{Apa yang saya kerjakan.} Ringkasan tindakan pada
          sesi ini.
    \item \textbf{Apa yang menghambat.} Kendala teknis atau konseptual
          yang muncul.
    \item \textbf{Bagaimana saya mengatasinya.} Cara mengatasi beserta
          alasan mengapa cara itu benar.
    \item \textbf{Apa yang saya pahami sekarang.} Satu pemahaman baru
          yang diperoleh, dinyatakan dengan kalimat sendiri.
\end{enumerate}

\section{Rubrik Penilaian Refleksi}

\begin{table}[H]
    \centering
    \caption{Rubrik penilaian jurnal refleksi.}
    \label{tab:app-rubrik-refleksi}
    \footnotesize
    \begin{tabular}{@{}p{3.0cm}p{3.4cm}p{3.4cm}p{3.4cm}@{}}
        \toprule
        \textbf{Kriteria} & \textbf{4 -- Sangat baik} &
        \textbf{3 -- Baik} & \textbf{2 -- Cukup} \\
        \midrule
        Kedalaman analisis & Menemukan penyebab akar masalah, bukan
        hanya gejalanya & Menjelaskan penyebab dengan cukup &
        Menyebutkan gejala tanpa analisis \\
        \addlinespace
        Kejujuran & Menuliskan kegagalan dan kesalahpahaman secara
        terbuka & Menuliskan sebagian kendala & Hanya menuliskan
        keberhasilan \\
        \addlinespace
        Keterhubungan & Mengaitkan pengalaman dengan konsep pada bab
        yang relevan & Mengaitkan secara umum & Tidak mengaitkan \\
        \addlinespace
        Rencana tindak lanjut & Menyusun langkah perbaikan yang
        konkret dan terukur & Menyusun langkah umum & Tidak
        menyusun rencana \\
        \bottomrule
    \end{tabular}
\end{table}

\section{Rubrik Penilaian Presentasi}

\begin{table}[H]
    \centering
    \caption{Rubrik penilaian presentasi proyek akhir.}
    \label{tab:app-rubrik-presentasi}
    \footnotesize
    \begin{tabular}{@{}p{3.0cm}p{3.4cm}p{3.4cm}p{3.4cm}@{}}
        \toprule
        \textbf{Kriteria} & \textbf{4 -- Sangat baik} &
        \textbf{3 -- Baik} & \textbf{2 -- Cukup} \\
        \midrule
        Kejelasan penyampaian & Alur runtut, waktu terjaga, istilah
        tepat & Alur cukup runtut & Alur melompat-lompat \\
        \addlinespace
        Penguasaan materi & Menjawab pertanyaan dengan menunjuk kode
        secara langsung & Menjawab sebagian besar pertanyaan &
        Menjawab dengan ragu \\
        \addlinespace
        Kualitas demonstrasi & Program berjalan lancar, kasus salah
        diperlihatkan & Program berjalan lancar & Program berjalan
        setelah dicoba ulang \\
        \bottomrule
    \end{tabular}
\end{table}
